\uuid{4cvr}
\exo7id{3564}
\titre{exo7 3564}
\auteur{quercia}
\organisation{exo7}
\datecreate{2010-03-10}
\isIndication{false}
\isCorrection{true}
\chapitre{Réduction d'endomorphisme, polynôme annulateur}
\sousChapitre{Polynôme annulateur}
\cpge{mpsi-mp2i.changements-bases.b, mp-mpi.reduction-endomorphismes.a, pc.complements-espaces-vectoriels.a, psi.complements-espaces-vectoriels.a, pt.complements-algebre-lineaire.a, mp-mpi.reduction-endomorphismes.c, pc.reduction.b, psi.reduction-endomorphismes.b, pt.reduction-endomorphismes.b, pc.complements-espaces-vectoriels.b, psi.complements-espaces-vectoriels.b, pt.complements-algebre-lineaire.b}
\module{Algèbre}
\niveau{L2}
\difficulte{}

\contenu{
\question{
Soient $A,B,C\in\mathcal{M}_n(\C)$ telles que $AC = CB$ et $\mathrm{rg}(C) = r$.
Montrer que $A$ et $B$ ont au moins $r$ valeurs propres communes.
}
\reponse{
On écrit $C = PJQ$ où $P,Q$ sont inversibles et $J$ est la matrice
canonique de rang~$r$. Alors $(P^{-1}AP)J = J(QBQ^{-1})$ donc $P^{-1}AP$
et $QBQ^{-1}$ sont triangulaires par blocs avec le même bloc diagonal
$r\times r$, ce qui prouve que $\chi_A$ et $\chi_B$ ont un facteur de
degré~$r$ en commun.
}
}
